% LaTeX summary snippet (to be \input into a larger document)

\section*{Complex Summary of 1285\_MS Research Paper}

\subsection*{English Summary}
\subsubsection*{Title: Teaching Precommitted Agents: Model-Free Policy Evaluation and Control in Quasi-Hyperbolic Discounted MDPs}
\textbf{Author:} S.R. Eshwar, Department of Computer Science and Automation, Indian Institute of Science, Bengaluru, India

\paragraph{Research Overview}
This work investigates time-inconsistent preferences via quasi-hyperbolic (QH) discounting in RL, moving beyond exponential discounting and time-consistent preferences to capture present-bias.

\paragraph{Core Problem Statement}
In real decision-making, preferences between two future rewards may reverse as time advances (present-bias). Modeling and learning under such time inconsistency requires new theory and algorithms.

\subsubsection*{Key Theoretical Contributions}
\begin{enumerate}
  \item \textbf{Optimal policy structure (one-step non-stationary):} The optimal QH policy decomposes into a first-step policy $\mu^*$ and a stationary continuation policy $\bar{\pi}^*$.
  \item \textbf{Mathematical framework:} Cumulative return
  \[
    G = r_0 + \sigma \sum_{t=1}^{\infty} \gamma^t r_t,\quad \sigma\in[0,1],\; \gamma\in[0,1)
  \]
  encapsulates present-bias $\sigma$ and exponential factor $\gamma$.
  \item \textbf{Agent types:} \emph{Naive} (replan without anticipating inconsistency), \emph{Sophisticated} (subgame-perfect planning), \emph{Precommitted} (fix a policy to maximize QH return).
\end{enumerate}

\subsubsection*{Algorithmic Innovations}
\begin{enumerate}
  \item \textbf{Model-free policy evaluation:} Two-time-scale stochastic approximation for QH evaluation with convergence guarantees.
  \item \textbf{QH Q-learning:} Separately estimate exponential and QH Q-functions; off-policy learning converges to the optimal one-step non-stationary pair $(\mu^*,\bar{\pi}^*)$.
\end{enumerate}

\subsubsection*{Experimental Validation: Inventory Control}
Parameters: $M=2$, $c=5$, $h=2$, $p=9$, $\Pr(D{=}0,1,2)=(0.2,0.3,0.5)$, $(\sigma,\gamma)=(0.3,0.9)$. Results: recovered optimal policy pairs by value iteration and QH Q-learning; empirical convergence matched theory.

\subsubsection*{Broader Implications and Applications}
\begin{itemize}
  \item Behavioral economics integration: formal RL tools for present-biased decisions.
  \item Applications: finance, healthcare adherence, inventory/energy/supply chains, human–AI interaction.
  \item Methodological advances: first model-free algorithms for QH in RL; scalable computational tools.
\end{itemize}

\subsubsection*{Technical Significance}
\begin{itemize}
  \item \textbf{Rigor:} Policy structure theorems; ODE-based convergence analysis.
  \item \textbf{Efficiency:} Reduce infinite-horizon non-stationarity to two-component policies; model-free and scalable.
\end{itemize}

\subsubsection*{Future Directions}
Extensions to POMDPs, multi-agent settings, continuous spaces; applications to dynamic pricing, robotics, recommendation; empirical validation and parameter estimation for present-bias.

\paragraph{Conclusion}
QH discounting yields tractable two-component optimal policies, enabling theoretically sound, practical RL for time-inconsistent preferences and unlocking realistic applications.

\bigskip
\subsection*{Polish Summary / Polskie Streszczenie}
\subsubsection*{Tytuł: Nauczanie Zaangażowanych Agentów: Wolna od Modeli Ocena Polityki i Kontrola w MDP z Quasi-Hiperbolicznym Dyskontowaniem}
\textbf{Autor:} S.R. Eshwar, Wydział Informatyki i Automatyki, Indyjski Instytut Nauki, Bengaluru, Indie

\paragraph{Przegląd Badań}
Praca bada preferencje niespójne w czasie poprzez quasi-hiperboliczne dyskontowanie, wykraczając poza klasyczne, wykładnicze dyskontowanie i spójność w czasie.

\paragraph{Główny Problem Badawczy}
Preferencje mogą się odwracać wraz z upływem czasu (skłonność do teraźniejszości), co wymaga nowych narzędzi teoretycznych i algorytmicznych w RL.

\subsubsection*{Kluczowe Wkłady Teoretyczne}
\begin{enumerate}
  \item \textbf{Struktura polityk:} Optimum ma formę jednostopniowej polityki niestacjonarnej: pierwszy krok $\mu^*$, dalej stacjonarna $\bar{\pi}^*$.
  \item \textbf{Rama matematyczna:}
  \(
    G = r_0 + \sigma \sum_{t=1}^{\infty} \gamma^t r_t,\; \sigma\in[0,1],\; \gamma\in[0,1)
  \).
  \item \textbf{Typy agentów:} naiwni, wyrafinowani, zaangażowani (precommitment).
\end{enumerate}

\subsubsection*{Innowacje Algorytmiczne}
\begin{itemize}
  \item Dwu-skalowa aproksymacja stochastyczna do oceny polityk QH z gwarancjami zbieżności.
  \item QH Q-learning: rozdzielenie estymacji Q i zbieżność do pary $(\mu^*,\bar{\pi}^*)$.
\end{itemize}

\subsubsection*{Walidacja Eksperymentalna (Zapas)}
Parametry: $M{=}2$, $c{=}5$, $h{=}2$, $p{=}9$, $\Pr(D)=(0.2,0.3,0.5)$, $(\sigma,\gamma)=(0.3,0.9)$. Wyniki: odzyskano optymalne pary polityk; zbieżność algorytmu potwierdzona empirycznie.

\subsubsection*{Implikacje i Zastosowania}
Integracja z ekonomią behawioralną; zastosowania w finansach, zdrowiu, zarządzaniu zasobami, interakcji człowiek–AI; narzędzia skalowalne.

\subsubsection*{Znaczenie Techniczne i Kierunki}
Rygor dowodów i efektywność obliczeń; rozszerzenia do POMDP, systemów wieloagentowych i przestrzeni ciągłych; walidacja empiryczna.

\paragraph{Podsumowanie}
QH upraszcza złożone polityki do eleganckich struktur dwuskładnikowych, czyniąc analizę i obliczenia wykonalnymi w realistycznych zadaniach RL z niespójnością czasową.
